\documentclass[a4paper,12pt]{article}

\usepackage[T2A]{fontenc}
\usepackage[utf8]{inputenc}
\usepackage[russian]{babel}
\usepackage{amsmath,amssymb}
\usepackage{PTSans}
\renewcommand{\familydefault}{\sfdefault}
\usepackage{icomma}
\usepackage[a4paper,left=20mm,right=20mm,top=20mm,bottom=22mm]{geometry}
\usepackage[nopatch=footnote]{microtype}
\usepackage{tikz}
\usetikzlibrary{calc}
\usetikzlibrary{arrows.meta,positioning,shapes.geometric}
\usepackage{tabularx,booktabs}
\usepackage{enumitem}
\usepackage{hyperref}
\usepackage{parskip}
\usepackage{fancyhdr}
\usepackage{setspace}
\usepackage{xcolor}

\definecolor{accent}{HTML}{008080}
\definecolor{darkaccent}{HTML}{004D4D}
\definecolor{textgray}{HTML}{333333}

\hypersetup{
    colorlinks=true,
    linkcolor=darkaccent,
    urlcolor=darkaccent
}

\onehalfspacing
\setlength{\parindent}{0pt}
\setlist[itemize]{
    leftmargin=1.6em,
    itemsep=0.25em,
    topsep=0.25em
}
\setlist[enumerate]{
    leftmargin=1.8em,
    itemsep=0.45em,
    topsep=0.35em
}

\renewcommand{\arraystretch}{1.35}

\pagestyle{fancy}
\fancyhf{}
\fancyhead[L]{\small\textcolor{textgray}{Стереометрия: аксиомы и следствия}}
\fancyhead[R]{\small\textcolor{textgray}{Ксения Клыкова}}
\fancyfoot[C]{\small\thepage}
\setlength{\headheight}{24.5pt}
\renewcommand{\headrulewidth}{0.3pt}
\renewcommand{\headrule}{%
    \hbox to\headwidth{%
        \color{accent}\leaders\hrule height \headrulewidth\hfill
    }%
}

\newcommand{\term}[1]{\textbf{\textcolor{darkaccent}{#1}}}

\tikzset{
    flowstart/.style={
        ellipse,
        draw=darkaccent,
        line width=0.9pt,
        align=center,
        minimum width=4.1cm,
        minimum height=1.0cm,
        inner sep=6pt
    },
    flowaction/.style={
        rounded corners=3pt,
        draw=accent,
        line width=0.9pt,
        align=center,
        text width=8.7cm,
        minimum height=1.0cm,
        inner sep=7pt
    },
    flowdecision/.style={
        diamond,
        aspect=2.2,
        draw=darkaccent,
        line width=0.9pt,
        align=center,
        text width=5.0cm,
        inner sep=3pt
    },
    flowarrow/.style={
        -{Stealth[length=3mm,width=2mm]},
        line width=0.9pt,
        draw=textgray
    }
}

\begin{document}

% =========================================================
% ТИТУЛЬНЫЙ ЛИСТ
% =========================================================

\begin{titlepage}
\thispagestyle{empty}
\centering

\vspace*{2.2cm}

{\Large\textcolor{darkaccent}{Чек-лист}\par}

\vspace{0.7cm}

{\huge\bfseries
Аксиомы стереометрии и способы задания плоскости
\par}

\vspace{1.1cm}

{\large
Следствия: прямая и точка, пересекающиеся и параллельные прямые
\par}

\vfill

{\large Ксения Клыкова\par}

\vspace{0.3cm}

{\normalsize
Подготовка к ЕГЭ по профильной математике
\par}

\vspace{0.8cm}

{\normalsize \today\par}

\vfill

{\normalsize
Лёвин Артём Александрович эксклюзивно для Ксении Клыковой
\par}

\vspace{1.1cm}

\begin{tikzpicture}[x=1cm,y=1cm]
    \draw[
        accent,
        line width=1pt
    ]
    (-7.2,0)
    .. controls (-4.8,0.55) and (-2.6,-0.45) .. (0,0)
    .. controls (2.6,0.45) and (4.8,-0.55) .. (7.2,0);
\end{tikzpicture}

\end{titlepage}

\setcounter{page}{1}

% =========================================================
% ОСНОВНОЙ ЧЕК-ЛИСТ
% =========================================================

\section*{1. Что нужно помнить}

\subsection*{1.1. Основные обозначения}

\begin{itemize}
    \item $A\in a$ --- точка $A$ лежит на прямой $a$.
    \item $A\in\alpha$ --- точка $A$ лежит в плоскости $\alpha$.
    \item $a\subset\alpha$ --- вся прямая $a$ лежит в плоскости $\alpha$.
    \item $a\cap b=\{O\}$ --- прямые $a$ и $b$ пересекаются в точке $O$.
    \item $(ABC)$ --- плоскость, проходящая через три точки
    $A,B,C$, если эти точки не лежат на одной прямой.
\end{itemize}

\subsection*{1.2. Три аксиомы стереометрии}

\begin{enumerate}
    \item[\textbf{A1.}]
    Через любые три точки, не лежащие на одной прямой,
    проходит \term{единственная плоскость}.

    \item[\textbf{A2.}]
    Если две точки прямой лежат в плоскости,
    то вся эта прямая лежит в плоскости.

    \item[\textbf{A3.}]
    Если две плоскости имеют общую точку,
    то они имеют общую прямую, на которой лежат
    все их общие точки.
\end{enumerate}

\textbf{Ключевое слово:} \term{единственная}.

Оно фиксирует сразу два факта:

\begin{itemize}
    \item требуемая плоскость существует;
    \item второй отличной плоскости с теми же определяющими
    объектами быть не может.
\end{itemize}

% =========================================================

\section*{2. Три главных способа однозначно задать плоскость}

\begin{table}[htbp]
\centering
\caption{Опорная карта способов задания плоскости}

\begin{tabularx}{\linewidth}{
    @{}
    p{0.28\linewidth}
    X
    X
    @{}
}
\toprule
\textbf{Определяющие объекты}
&
\textbf{Условие}
&
\textbf{Основание}
\\
\midrule

Три точки $A,B,C$
&
Не лежат на одной прямой
&
Аксиома A1
\\

Прямая $a$ и точка $M$
&
$M\notin a$
&
Следствие 1
\\

Две прямые $a,b$
&
$a\cap b\ne\emptyset$
&
Следствие 2
\\

Две прямые $a,b$
&
$a\parallel b$
&
Следствие 3
\\

\bottomrule
\end{tabularx}
\end{table}

% =========================================================

\subsection*{2.1. Следствие 1: прямая и точка вне неё}

\term{Формулировка.}

Через прямую $a$ и точку $M\notin a$
проходит единственная плоскость.

\textbf{Доказательство.}

Выберем на прямой $a$ две различные точки $A$ и $B$:
\[
A\in a,
\qquad
B\in a,
\qquad
A\ne B.
\]

Поскольку
\[
M\notin a,
\]
точки $A,B,M$ не лежат на одной прямой.

По A1 через них проходит единственная плоскость $\alpha$.

Имеем
\[
A,B\in\alpha
\qquad\text{и}\qquad
A,B\in a.
\]

По A2:
\[
a\subset\alpha.
\]

Таким образом, плоскость $\alpha$ содержит одновременно
прямую $a$ и точку $M$.

Единственность также следует из A1:
любая плоскость, содержащая $a$ и $M$,
содержит точки $A,B,M$.

% =========================================================

\subsection*{2.2. Следствие 2: две пересекающиеся прямые}

\term{Формулировка.}

Через две пересекающиеся прямые проходит
единственная плоскость.

Пусть
\[
a\cap b=\{O\}.
\]

Выберем
\[
A\in a,
\qquad
A\ne O,
\]
и
\[
B\in b,
\qquad
B\ne O.
\]

Точки $A,O,B$ не лежат на одной прямой.

Следовательно, по A1 они задают единственную плоскость
\[
\alpha=(AOB).
\]

Так как
\[
A,O\in\alpha
\qquad\text{и}\qquad
A,O\in a,
\]
по A2 получаем
\[
a\subset\alpha.
\]

Аналогично:
\[
B,O\in\alpha,
\qquad
B,O\in b,
\]
поэтому
\[
b\subset\alpha.
\]

Пусть некоторая плоскость $\beta$ также содержит
обе прямые $a$ и $b$.

Тогда
\[
A,O,B\in\beta.
\]

По единственности плоскости через три точки,
не лежащие на одной прямой,
\[
\boxed{\alpha=\beta}.
\]

\textbf{Почему важно требовать}
\[
A\ne O,
\qquad
B\ne O?
\]

Для применения A2 требуются две различные точки одной прямой.

% =========================================================

\subsection*{2.3. Следствие 3: две параллельные прямые}

\term{Определение.}

Две прямые называются параллельными,
если они лежат в одной плоскости
и не имеют общих точек.

\term{Формулировка.}

Через две параллельные прямые проходит
единственная плоскость.

\textbf{Доказательство.}

Пусть
\[
a\parallel b.
\]

По определению параллельных прямых существует плоскость,
содержащая $a$ и $b$.

Выберем точку
\[
M\in b.
\]

Поскольку параллельные прямые общих точек не имеют,
\[
M\notin a.
\]

По следствию 1 через прямую $a$ и точку $M$
проходит единственная плоскость $\alpha$.

Любая плоскость, содержащая одновременно $a$ и $b$,
обязательно содержит:

\[
a
\qquad\text{и}\qquad
M\in b.
\]

Следовательно, по единственности из следствия 1
такая плоскость совпадает с $\alpha$.

Итак, две параллельные прямые задают
\[
\boxed{\text{единственную плоскость}}.
\]

% =========================================================

\section*{3. Как использовать единственность в задачах}

Типичная задача просит доказать совпадение
двух плоскостей.

Рабочий принцип:

\begin{enumerate}
    \item найти объекты, принадлежащие обеим плоскостям;
    \item проверить, что эти объекты однозначно задают плоскость;
    \item сослаться на соответствующую аксиому или следствие;
    \item сделать вывод о совпадении плоскостей.
\end{enumerate}

\begin{itemize}
    \item Если обе плоскости содержат две пересекающиеся прямые,
    используйте следствие 2.

    \item Если обе плоскости содержат две параллельные прямые,
    используйте следствие 3.

    \item Если обе плоскости содержат одну прямую
    и точку вне неё, используйте следствие 1.

    \item Если обе плоскости содержат три точки,
    не лежащие на одной прямой, используйте A1.
\end{itemize}

\textbf{Шаблон вывода:}

\[
\text{одни и те же определяющие объекты}
\quad\Longrightarrow\quad
\text{единственная плоскость}
\quad\Longrightarrow\quad
\alpha=\beta.
\]

% =========================================================
% БЛОК-СХЕМА НА ОТДЕЛЬНОЙ СТРАНИЦЕ
% =========================================================

\clearpage
\thispagestyle{empty}

\begin{center}
{\Large\bfseries\textcolor{darkaccent}{
Алгоритм: доказать совпадение двух плоскостей
}}
\end{center}

\vspace{0.5cm}

\begin{center}
\begin{tikzpicture}[node distance=1.15cm]

\node[flowstart] (start)
{Нужно доказать $\alpha=\beta$};

\node[
    flowaction,
    below=of start
] (objects)
{Найдите объекты, которые лежат одновременно
в $\alpha$ и $\beta$};

\node[
    flowdecision,
    below=1.35cm of objects
] (direct)
{Есть стандартный набор,
однозначно задающий плоскость?};

\node[
    flowaction,
    below=1.35cm of direct
] (choose)
{Выберите основание:
A1; следствие 1; следствие 2; следствие 3};

\node[
    flowaction,
    below=of choose
] (unique)
{Сошлитесь на единственность
плоскости для выбранного набора};

\node[
    flowstart,
    below=of unique
] (finish)
{Сделайте вывод $\alpha=\beta$};

\node[
    flowaction,
    right=2.0cm of direct,
    text width=4.3cm
] (extract)
{Выделите три неколлинеарные точки
либо прямую и точку вне неё;
при необходимости примените A2};

\draw[flowarrow]
(start) -- (objects);

\draw[flowarrow]
(objects) -- (direct);

\draw[flowarrow]
(direct)
--
node[left,font=\small]{да}
(choose);

\draw[flowarrow]
(choose) -- (unique);

\draw[flowarrow]
(unique) -- (finish);

\draw[flowarrow]
(direct.east)
--
node[above,font=\small]{нет}
(extract.west);

\draw[flowarrow]
(extract.south)
|-
($(choose.east)+(0,0.35)$);

\end{tikzpicture}
\end{center}

\clearpage

% =========================================================

\section*{4. Связь с треугольником}

Если $AB$ и $AC$ --- стороны невырожденного
треугольника $ABC$, то прямые $AB$ и $AC$
пересекаются в точке $A$.

По следствию 2 они задают единственную плоскость.

Эта плоскость является плоскостью треугольника $ABC$.

Тот же результат можно получить через A1:

\[
A,\ B,\ C
\]
не лежат на одной прямой, поэтому задают единственную
плоскость
\[
(ABC).
\]

Итак,
\[
\boxed{
AB\ \text{и}\ AC
\text{ однозначно задают плоскость треугольника }ABC
}.
\]

% =========================================================

\section*{5. Когда две прямые плоскость не задают}

В пространстве две различные прямые могут:

\begin{itemize}
    \item пересекаться;
    \item быть параллельными;
    \item быть \term{скрещивающимися}.
\end{itemize}

\term{Скрещивающиеся прямые} ---
это прямые, которые не лежат в одной плоскости.

Следовательно, общей плоскости,
содержащей обе скрещивающиеся прямые целиком,
не существует.

Важно:

\[
a\cap b=\emptyset
\]
ещё недостаточно для вывода
\[
a\parallel b.
\]

Для параллельности требуется совместная принадлежность
одной плоскости.

% =========================================================

\section*{6. Типичные ошибки}

\begin{enumerate}
    \item
    Использовать формулировку
    «через объекты проходит одна плоскость»
    без указания единственности.

    Корректная ключевая формулировка:
    \[
    \boxed{\text{проходит единственная плоскость}}.
    \]

    \item
    Делать вывод
    \[
    a\subset\alpha
    \]
    только по одной известной точке прямой в плоскости.

    Для A2 нужны две различные точки:
    \[
    A,B\in a,
    \qquad
    A,B\in\alpha,
    \qquad
    A\ne B.
    \]

    \item
    Считать любые две непересекающиеся прямые
    параллельными.

    В пространстве возможны скрещивающиеся прямые.

    \item
    При доказательстве следствия 2 забывать,
    что точки $A,O,B$ должны не лежать на одной прямой.

    \item
    При доказательстве
    \[
    \alpha=\beta
    \]
    перечислять множество фактов без указания,
    какой набор объектов однозначно задаёт плоскость.

    \item
    Путать принадлежность точки
    \[
    A\in\alpha
    \]
    и принадлежность прямой
    \[
    a\subset\alpha.
    \]

    \item
    Путать слова «одна» и «единственная»
    в доказательствах на однозначность конструкции.
\end{enumerate}

% =========================================================

\section*{7. Мини-памятка}

\begin{table}[htbp]
\centering
\caption{Что распознать в условии}

\begin{tabularx}{\linewidth}{@{}X X@{}}
\toprule

\textbf{В условии видно}
&
\textbf{Сразу вспоминайте}
\\
\midrule

$A,B,C$ не лежат на одной прямой
&
A1:
три точки $\Rightarrow$ единственная плоскость
\\

$a$ и $M\notin a$
&
Следствие 1
\\

$a\cap b\ne\emptyset$
&
Следствие 2
\\

$a\parallel b$
&
Следствие 3
\\

Две плоскости содержат
один и тот же определяющий набор
&
Используйте единственность и докажите
совпадение плоскостей
\\

Две прямые не пересекаются
&
Проверьте, лежат ли они в одной плоскости
\\

\bottomrule
\end{tabularx}
\end{table}

% =========================================================
% ТРЕНИРОВОЧНЫЙ БЛОК
% =========================================================

\clearpage

\section*{8. Тренировочный блок}

\textbf{Уровень: от среднего к сложному.}

Решения оформляйте кратко.
В каждом доказательстве обязательно указывайте
используемую аксиому или следствие.

\subsection*{Средний уровень}

\begin{enumerate}

    \item
    Прямые $a$ и $b$ пересекаются в точке $O$.
    Сколько плоскостей проходит одновременно через обе прямые?
    Укажите основание.

    \item
    Параллельные прямые $m$ и $n$ лежат
    и в плоскости $\alpha$, и в плоскости $\beta$.
    Докажите, что
    \[
    \alpha=\beta.
    \]

    \item
    Прямые $AB$ и $AC$ являются сторонами
    невырожденного треугольника $ABC$.
    Объясните, почему они однозначно задают
    плоскость треугольника.

\end{enumerate}

\subsection*{Выше среднего}

\begin{enumerate}[resume]

    \item
    Пусть
    \[
    a\cap b=\{O\},
    \qquad
    A\in a,
    \qquad
    A\ne O,
    \qquad
    B\in b,
    \qquad
    B\ne O.
    \]

    Докажите, что плоскость,
    заданная прямыми $a$ и $b$,
    совпадает с плоскостью
    \[
    (AOB).
    \]

    \item
    Пусть
    \[
    a\parallel b,
    \qquad
    M\in b.
    \]

    Докажите, что плоскость,
    заданная прямыми $a$ и $b$,
    совпадает с плоскостью,
    заданной прямой $a$ и точкой $M$.

    \item
    Две плоскости $\alpha$ и $\beta$
    содержат пересекающиеся прямые $p$ и $q$.

    Докажите совпадение плоскостей.

    \item
    На прямой $a$ выбраны различные точки $A$ и $B$,
    точка $C$ лежит на прямой $b$, причём
    \[
    a\parallel b.
    \]

    Докажите, что точки $A,B,C$
    не лежат на одной прямой
    и что плоскость $(ABC)$ совпадает
    с плоскостью, заданной $a$ и $b$.

    \item
    Прямая $l$ и точка
    \[
    M\notin l
    \]
    лежат и в плоскости $\alpha$,
    и в плоскости $\beta$.

    Докажите:
    \[
    \alpha=\beta.
    \]

    \item
    В плоскости $\alpha$ лежат параллельные
    прямые $a$ и $b$.

    В плоскости $\beta$ лежит прямая $a$
    и точка
    \[
    M\in b.
    \]

    Докажите:
    \[
    \alpha=\beta.
    \]

\end{enumerate}

\subsection*{Сложная задача}

\begin{enumerate}[resume]

    \item
    Прямые $a$ и $b$ скрещиваются.

    На прямой $b$ выбраны различные точки
    \[
    M
    \qquad\text{и}\qquad
    N.
    \]

    Пусть $\alpha$ --- плоскость,
    заданная прямой $a$ и точкой $M$,
    а $\beta$ --- плоскость,
    заданная прямой $a$ и точкой $N$.

    Докажите:
    \[
    \alpha\ne\beta.
    \]

\end{enumerate}

% =========================================================
% ОТВЕТЫ
% =========================================================

\clearpage

\section*{9. Ответы к тренировочному блоку}

\begin{table}[htbp]
\centering
\caption{Краткие ответы и ключевые основания}

\begin{tabularx}{\linewidth}{
    @{}
    p{0.07\linewidth}
    X
    @{}
}
\toprule
\textbf{№}
&
\textbf{Ответ / ключевая идея}
\\
\midrule

1
&
Единственная плоскость; следствие 2.
\\

2
&
\[
\alpha=\beta.
\]
Две параллельные прямые задают
единственную плоскость; следствие 3.
\\

3
&
\[
AB\cap AC=\{A\}.
\]
Две пересекающиеся прямые задают
единственную плоскость; следствие 2.
\\

4
&
Плоскость, заданная $a$ и $b$,
совпадает с $(AOB)$:
обе содержат три неколлинеарные точки
$A,O,B$.
Используются A1 и A2.
\\

5
&
Плоскости совпадают.
Имеем
\[
M\in b,
\qquad
M\notin a.
\]
Прямая $a$ и точка $M$
задают единственную плоскость;
следствие 1.
\\

6
&
\[
\alpha=\beta.
\]
Основание: следствие 2.
\\

7
&
Из
\[
a\parallel b
\]
следует
\[
C\notin a.
\]
Поэтому $A,B,C$ не лежат на одной прямой.
Плоскость, заданная $a,b$,
содержит $A,B,C$;
по A1 она совпадает с $(ABC)$.
\\

8
&
\[
\alpha=\beta.
\]
Основание: следствие 1.
\\

9
&
Плоскость $\alpha$ содержит $a$ и $M$,
поскольку
\[
M\in b\subset\alpha.
\]
Плоскость $\beta$ также содержит $a$ и $M$.
При этом
\[
M\notin a.
\]
По следствию 1:
\[
\alpha=\beta.
\]
\\

10
&
\[
\alpha\ne\beta.
\]
При предположении
\[
\alpha=\beta
\]
одна плоскость содержала бы $a$, $M$ и $N$.
Так как
\[
M,N\in b,
\qquad
M\ne N,
\]
по A2 получили бы
\[
b\subset\alpha.
\]
Тогда $a$ и $b$ лежали бы в одной плоскости,
что противоречит условию об их скрещивании.
\\

\bottomrule
\end{tabularx}
\end{table}

\end{document}